\section{Method}
\label{sec:method}

\subsection{Models and Connectors}
\label{sec:models}

We evaluate token pruning across six VLA models spanning three vision connector families (Table~\ref{tab:model_properties}). These connectors differ in whether they process ViT patches independently or create cross-token dependencies, a distinction we formalize as \emph{coupling}.

\textbf{Linear projection} applies a single matrix multiplication per patch. $\pi_0$~\cite{black2024pi0} projects 256 patches through a linear layer; $\pi_0.5$~\cite{black2025pi05} uses a $\times\!\sqrt{d}$-scaled linear projection across 768 tokens from three cameras. Both produce tokens with near-uniform L2 norms (CV $< 0.06$).

\textbf{MLP projector} maps each patch through a multi-layer perceptron. OpenVLA~\cite{kim2024openvla} and OFT~\cite{kim2025openoft} both use three-layer MLPs to produce 256 and 512 tokens respectively. Each patch is processed independently, and the resulting norms exhibit a clear importance gradient (CV $\approx 0.22$).

\textbf{Token-entangling connectors} introduce cross-token dependencies through post-projection processing. X-VLA~\cite{zheng2025xvla} feeds linearly projected tokens through a 12-layer BART encoder~\cite{lewis2020bart} to produce 50 tokens with learned cross-token dependencies via self-attention. SmolVLA~\cite{shukor2025smolvla} applies pixel-shuffle $4{\times}4$ folding that concatenates 16 contiguous patches into the channel dimension of each output token, for a total of 128 tokens.

We classify coupling as \emph{independent} when each output token depends on a single input patch, and \emph{entangling} when output tokens encode information from multiple patches. Four models ($\pi_0.5$, X-VLA, OpenVLA, SmolVLA) are evaluated in closed-loop robotic control on LIBERO~\cite{liu2023libero}; $\pi_0$ and OFT are analyzed via offline metrics only.

\begin{table}[t]
\centering
\small
\caption{VLA models evaluated in this study. Coupling indicates whether the connector processes patches independently or creates cross-token dependencies. $N$ = number of vision tokens input to the LLM backbone. Models are grouped by evaluation protocol and ordered by CV within each group.}
\label{tab:model_properties}
\begin{tabular}{lrlcrll}
\toprule
Model & Params & Connector & $N$ & CV & Coupling & Eval \\
\midrule
$\pi_0.5$ & 4B & Linear ($\times\!\sqrt{d}$) & 768 & 0.055 & Indep.\ & Closed-loop \\
X-VLA & 880M & Linear + BART & 50 & 0.151 & Entangling & Closed-loop \\
OpenVLA & 7B & 3-layer MLP & 256 & 0.215 & Indep.\ & Closed-loop \\
SmolVLA & 450M & Pixel-shuffle $4{\times}4$ & 128 & 0.314 & Entangling & Closed-loop \\
\midrule
$\pi_0$ & 3.3B & Linear & 256 & 0.038 & Indep.\ & Offline \\
OFT & 7B & 3-layer MLP & 512 & 0.219 & Indep.\ & Offline \\
\bottomrule
\end{tabular}
\end{table}

\subsection{Pruning Protocol}
\label{sec:pruning_protocol}

We apply post-connector token pruning at retention ratio $k \in \{0.3, 0.5, 0.7, 0.85, 0.9, 0.95\}$, where $k$ denotes the fraction of $N$ vision tokens retained. Given $N$ post-connector tokens, pruning produces $\lfloor k \cdot N \rfloor$ tokens that are passed to the LLM backbone.

Two strategies bracket expected performance:

\textbf{L2 norm selection} ranks tokens by their L2 norm and retains the top-$k$ fraction. This importance-based strategy provides a favorable test: if a genuine importance hierarchy exists among tokens, L2 selection should outperform random baselines.

\textbf{Random subsampling} selects $\lfloor k \cdot N \rfloor$ tokens uniformly at random. Because it ignores token content entirely, random selection provides a conservative lower bound.

Both strategies physically remove unselected tokens and re-index position IDs to close gaps in the sequence. Pruned tokens are discarded, not replaced with zero vectors or mask tokens.

\subsection{Compressibility Diagnostic: Importance CV}
\label{sec:cv_diagnostic}

We quantify the heterogeneity of post-connector token representations using the coefficient of variation (CV) of their L2 norms:
\begin{equation}
\text{CV} = \frac{\sigma\!\left(\{\|t_i\|_2\}_{i=1}^{N}\right)}{\mu\!\left(\{\|t_i\|_2\}_{i=1}^{N}\right)}
\label{eq:cv}
\end{equation}
where $t_i$ denotes the $i$-th post-connector vision token, and $\sigma$, $\mu$ are the standard deviation and mean computed over all $N$ token norms from a single observation.

CV is a single-pass diagnostic: it requires one forward pass through the vision encoder and connector, with no pruning or downstream evaluation needed. The metric captures a continuous spectrum across connector designs. At one end, $\pi_0$ (CV $= 0.038$) produces near-uniform norms, indicating that the linear connector distributes information evenly across patches. At the other end, SmolVLA (CV $= 0.314$) exhibits high norm heterogeneity from its pixel-shuffle operation. The remaining models span the interval between these extremes (Table~\ref{tab:model_properties}).

CV correlates with both pruning tolerance and sensitivity to pruning method choice. Models with low CV exhibit near-uniform tokens where the choice of \emph{which} tokens to remove matters more than \emph{how many} are removed. Models with high CV from entangling connectors exhibit catastrophic failure under any token removal strategy. We report success rate (SR) as the primary closed-loop evaluation metric throughout.

\subsection{Vision Token Re-encoder (VTR)}
\label{sec:vtr}

Hard pruning fails for token-entangling connectors because removing any token destroys spatially distributed information that cannot be recovered from the remaining tokens. Vision Token Re-encoder (VTR) addresses this limitation through learned re-encoding rather than hard deletion.

VTR is a lightweight cross-attention module with approximately 22M trainable parameters (roughly 5\% of SmolVLA's 450M total). The module takes the full set of $N$ source tokens from the vision connector as keys and values, and uses $M$ learned queries to attend over all $N$ tokens via multi-head cross-attention. Each query aggregates information from the entire spatial field. The resulting $M$ output representations are passed to the LLM backbone as compressed vision tokens. The output dimensionality $M \in \{16, 32, 64, 96, 128\}$ controls the compression ratio; $M = N = 128$ serves as a passthrough control that re-encodes without reducing the token count.

We also evaluate a parameter-matched MLP adapter control (approximately 22M parameters). This variant applies average pooling over groups of source tokens followed by a two-layer MLP, without cross-attention. Comparing VTR against this control isolates whether the performance gain stems from the cross-attention mechanism specifically or from the act of re-encoding entangled tokens into a new representation.
